\section{Referencias}

Las siguientes fuentes técnicas, normativas aeroespaciales y publicaciones académicas sustentan el marco metodológico y de diseño del presente informe, citadas conforme al formato estándar del \textit{American Institute of Aeronautics and Astronautics} (AIAA):

\begin{enumerate}[label={[\arabic*]}, leftmargin=2em]
    \item Concurso Colombiano de Cohetería Deportiva (CCCD), ``Documento de Diseño, Construcción y Pruebas (DCP 2026),'' Asociación Colombiana de Cohetería Experimental y Deportiva, Bogotá, Colombia, 2026.
    \item Concurso Colombiano de Cohetería Deportiva (CCCD), ``Norma de Seguridad en Cohetería Deportiva (NSRT),'' Sección 4: Estabilidad y Criterios de Aceptación Técnica, Bogotá, Colombia, 2026.
    \item Concurso Colombiano de Cohetería Deportiva (CCCD), ``Guía Oficial de Categoría Intermedios --- Motores Clase F,'' Medellín, Colombia, 2026.
    \item Concurso Colombiano de Cohetería Deportiva (CCCD), ``Guía Técnica de Carga Útil para Categoría Intermedios 2026,'' Medellín, Colombia, 2026.
    \item National Association of Rocketry (NAR), ``High Power Rocket Safety Code,'' Technical Services Committee, Marion, IA, EE.UU., 2023.
    \item Tripoli Rocketry Association (TRA), ``Safety Code and Range Safety Officer Guidelines,'' Valparaiso, IN, EE.UU., 2022.
    \item Barrowman, J. S., ``The Theoretical Prediction of the Center of Pressure,'' Centuri Engineering Co., TIR-33, Phoenix, AZ, 1966.
    \item Box, S., Bishop, C. M., y Hunt, H. E. M., ``Estimating the Dynamic Response and Stability of Model and High-Power Rockets,'' \textit{Journal of Aerospace Engineering}, Vol. 225, No. 3, 2011, págs. 305--318.
    \item Niskanen, S., ``OpenRocket Technical Documentation: Open Source Model Rocket Simulation Software,'' Helsinki University of Technology, Espoo, Finlandia, 2013.
    \item Estes Industries, ``Estes Model Rocket Engine Technical Data and Thrust Curves: F15 Series,'' Penrose, CO, EE.UU., 2021.
    \item Stine, G. H., y Stine, B., \textit{Handbook of Model Rocketry}, 7ª ed., John Wiley \& Sons, Hoboken, NJ, 2004.
    \item American Institute of Aeronautics and Astronautics (AIAA), ``AIAA Standards for Experimental Sounding Rockets and Suborbital Vehicles,'' Reston, VA, 2020.
\end{enumerate}
